弱-* 收敛概念的实用性在很大程度上基于以下结果.

\begin{Theorem}{弱紧性, Alaoglu}
\label{thm:alaoglu-weak-star-subsequence}
设 $X$ 为可分 Banach 空间, $f_n$ 为 $X^*$ 中的有界序列. 那么 $f_n$ 有一个弱-* 收敛子序列.
\end{Theorem}

\begin{Proof}
设 $x_k$, $k\in\mathbb{N}$ 为 $X$ 中的稠密序列. 首先注意, 若对每个 $k\in\mathbb{N}$, $f_n(x_k)$ 收敛, 则对每个 $x\in X$, $f_n(x)$ 收敛. 为说明这一点, 设 $M$ 为 $\|f_n\|$ 的一个上界. 对给定的 $\varepsilon>0$ 及 $x\in X$, 可以选取 $k$, 使 $\|x-x_k\|\leq\varepsilon/(3M)$. 于是有
\begin{equation}\label{eq:alaoglu-dense-set-cauchy-estimate}
\begin{aligned}
|f_n(x)-f_m(x)|&\leq|f_n(x)-f_n(x_k)|+|f_n(x_k)-f_m(x_k)|\\
&\qquad+|f_m(x_k)-f_m(x)|\\
&\leq\dfrac{2}{3}\varepsilon+|f_n(x_k)-f_m(x_k)|,
\end{aligned}
\end{equation}
当 $m$ 和 $n$ 充分大时, 右端最后一项小于 $\varepsilon/3$. 因而对每个 $x\in X$, $f_n(x)$ 收敛. 证明的其余部分是标准的对角线论证. 由于 $f_n(x_1)$ 有界, 可以抽取子序列 $f_{n1}$, 使 $f_{n1}(x_1)$ 收敛. 从这一子序列中再抽取子序列 $f_{n2}$, 使 $f_{n2}(x_2)$ 收敛. 如此归纳进行. 对角线序列 $f_{nn}$ 具有对每个 $k$, $f_{nn}(x_k)$ 都收敛的性质.
\end{Proof}

\begin{Remark}
\label{rem:weak-compactness-approximate-pde-solutions}
在该定理的典型应用中, $f_n$ 为某个 PDE 的近似解. 虽然可能难以证明该序列收敛, 但往往可以得到该序列的界, 从而得到一个弱收敛子序列. 然后证明这一子序列的极限为原 PDE 的解.
\end{Remark}
